% Sustentación: 8 diapositivas, 10 minutos, 5 expositores. Compilar: pdflatex diapositivas.tex (dos veces)
\documentclass[aspectratio=169,11pt]{beamer}
\usepackage[utf8]{inputenc}
\usepackage[T1]{fontenc}
\usepackage[spanish,es-noshorthands,es-tabla]{babel}
\usepackage{lmodern}
\usepackage{csquotes}
\usepackage{amssymb}
\usepackage{booktabs}
\usepackage{tikz}
\usetikzlibrary{arrows.meta,positioning}
\usepackage{pgfplots}
\pgfplotsset{compat=1.17}

% Mismos colores del informe (tinta) y de la app (rojo)
\definecolor{tinta}{HTML}{1F2A44}
\definecolor{rojo}{HTML}{C8102E}
\definecolor{suave}{HTML}{F3F4F7}
\usetheme{default}
\setbeamercolor{frametitle}{fg=tinta}
\setbeamercolor{title}{fg=tinta}
\setbeamercolor{structure}{fg=rojo}
\setbeamercolor{normal text}{fg=black!85}
\setbeamercolor{block title}{bg=tinta,fg=white}
\setbeamercolor{block body}{bg=suave}
\setbeamerfont{frametitle}{series=\bfseries,size=\Large}
\setbeamertemplate{navigation symbols}{}
\setbeamertemplate{itemize items}{\color{rojo}\small$\blacktriangleright$}
\setbeamertemplate{blocks}[rounded]
\setbeamertemplate{frametitle}{%
  \vspace{0.35cm}\insertframetitle\par
  \vspace{-0.15cm}{\color{rojo}\rule{\linewidth}{1.2pt}}}
% Quién expone cada diapositiva, en el pie
\newcommand{\expone}{}
\setbeamertemplate{footline}{%
  \hspace{0.4cm}{\color{gray}\scriptsize Expone: \textbf{\expone}}%
  \hfill{\color{gray}\scriptsize\insertframenumber/\inserttotalframenumber}\hspace{0.4cm}\vspace{0.25cm}}

\title{Enrutamiento de modelos de IA para el estudio universitario}
\subtitle{Un modelo de decisión con grafos, análisis de redes y una red neuronal}
\institute{Universidad del Norte \textbullet{} Matemáticas Discretas (MAT 4021) \textbullet{} Prof. MSc. Iván Gómez}
\date{16 de noviembre de 2026}

\begin{document}

% ================================================================ 1. Portada
{\setbeamertemplate{footline}{}
\begin{frame}
\vfill
{\color{rojo}\rule{\linewidth}{1.5pt}}\\[0.4cm]
{\LARGE\bfseries\color{tinta}\inserttitle}\\[0.25cm]
{\large\insertsubtitle}\\[0.3cm]
{\color{rojo}\rule{\linewidth}{1.5pt}}\\[0.6cm]
\small José Angulo \textbullet{} Sebastián Ibañez \textbullet{} Andrés Suarez \textbullet{} Samuel Ramírez \textbullet{} Manuel Hormechea\\[0.2cm]
{\color{gray}\insertinstitute\\ \insertdate}
\vfill
\end{frame}}

% ================================================================ 2. José: problema e impacto
\renewcommand{\expone}{José Angulo}
\begin{frame}{El problema: ¿qué modelo elegir al iniciar el chat?}
\small
\begin{columns}[T]
\column{0.5\textwidth}
\begin{itemize}
  \item En Colombia, el 84\,\% de los estudiantes encuestados usa IA generativa.
  \item Las universidades pasan a \textbf{proveerla}: CSU anunció acceso para más de 460\,000 estudiantes.
  \item Decenas de modelos con contexto, cupo y precios hasta 100 veces distintos.
\end{itemize}
\column{0.48\textwidth}
\begin{block}{Elegir mal cuesta}
\begin{itemize}
  \item \textbf{Costo:} el más potente para todo incrementa el gasto.
  \item \textbf{Confiabilidad:} elegir solo por precio puede afectar la calidad.
  \item \textbf{Continuidad:} si se satura, se corta la sesión.
\end{itemize}
\end{block}
\end{columns}
\vspace{0.25cm}
\begin{center}
\colorbox{suave}{\parbox{0.92\textwidth}{\centering Objetivo: al iniciar cada chat, elegir \textbf{\color{rojo}el modelo más barato que cumple las condiciones}.}}
\end{center}
{\footnotesize\textbf{Impacto:} afecta a instituciones y estudiantes de cualquier carrera \textbullet{} hay evidencia (precios de OpenRouter, puntajes GPQA) \textbullet{} se mide antes y después \textbullet{} se modela con la teoría del curso.}
\par\vspace{0.1cm}{\tiny\color{gray} Fuentes: Hernández (2024); California State University (2025); OpenRouter (2026). Referencias completas en el informe.}
\end{frame}

% ================================================================ 3. Sebastián: flujo, red neuronal y requisitos
\renewcommand{\expone}{Sebastián Ibañez}
\begin{frame}{Selección inicial: red neuronal y requisitos}
\begin{center}\scriptsize
Inicio $\to$ \textbf{\color{rojo}red neuronal} $\to$ requisitos $\to$ \textbf{\color{rojo}grafo bipartito} $\to$ \textbf{\color{rojo}red de sustitución} $\to$ \textbf{\color{rojo}ventana (inducción)} $\to$ API
\end{center}
\vspace{-0.3cm}
\vspace{0.2cm}
\begin{columns}[T]
\column{0.52\textwidth}
\textbf{Perceptrón multicapa} (hecho desde cero)
\[
h=\tanh(W_1x+b_1),\qquad N=\sigma(w_2^{\top}h+b_2)
\]
\begin{itemize}
  \item $x$: palabras, raíces de 5 letras y símbolos.
  \item 8 neuronas ocultas, retropropagación.
  \item 84 preguntas etiquetadas, varias materias.
  \item $N = P(\text{precisión en la entrada inicial})$.
\end{itemize}
\column{0.46\textwidth}
\begin{block}{Un modelo es apto si}
\begin{enumerate}
  \item está \textbf{disponible} en el catálogo conectado;
  \item su \textbf{contexto} alcanza: $\mathrm{ctx}\ge 1{,}5\cdot\text{tokens}$;
  \item su \textbf{puntaje} GPQA es $\ge 0{,}80$ si hay precisión ($N\ge0{,}5$ o la pide el estudiante); si no, $\ge 0{,}60$.
\end{enumerate}
\end{block}
{\scriptsize GPQA es una aproximación, no una garantía (Rein et al., 2023). En cada envío: contexto con margen $1{,}2$ y comprobación del cupo.}
\end{columns}
\end{frame}

% ================================================================ 4. Andrés: grafo bipartito
\renewcommand{\expone}{Andrés Suarez}
\begin{frame}{Grafo bipartito: selección inicial del modelo}
\begin{columns}[c]
\column{0.5\textwidth}
\begin{tikzpicture}[every node/.style={font=\small},
  t/.style={draw=tinta, circle, fill=tinta!12, minimum size=0.8cm},
  m/.style={draw, rounded corners, fill=suave, minimum width=3.9cm, minimum height=0.6cm, anchor=west}]
\node[t] (t1) at (0,1.5) {E}; \node[t] (t2) at (0,0) {P}; \node[t] (t3) at (0,-1.5) {R};
\node[m] (m1) at (2.2,2.1) {Gemma 4 26B \hfill \$0,08};
\node[m] (m2) at (2.2,0.7) {gpt-oss-120b \hfill \$0,15};
\node[m] (m3) at (2.2,-0.7) {Gemini 3.5 Flash \hfill \$1,50};
\node[m, dashed, text=gray] (m4) at (2.2,-2.1) {Llama 3.3 70B \hfill GPQA 0,50};
\foreach \t in {t1,t3} \foreach \m in {m1,m2,m3} \draw[gray!40] (\t.east) -- (\m.west);
\draw (t2.east) -- (m2.west); \draw (t2.east) -- (m3.west);
\draw[ultra thick, rojo] (t2.east) -- (m1.west);
\end{tikzpicture}
\column{0.5\textwidth}
\begin{itemize}
  \item Nodos: tareas $T$ y modelos $M$. Arista $\{t,m\}$ solo si $m$ es apto.
  \item Dos pesos por arista:
\end{itemize}
\[
\mathrm{costo}(t,m)=\frac{\text{tok}_{\text{ent}}\,p_{\text{ent}}+s_t\,p_{\text{sal}}}{10^6}
\]
\[
w(t,m)=\alpha_t\,\bar c(m)+\beta_t\,\mathrm{GPQA}+\gamma_t\,\mathrm{vel}
\]
{\scriptsize $\bar c(m)=\min(\mathrm{ctx}(m)/2^{20},1)$: contexto normalizado.}
\begin{block}{Regla}
$m^*=\arg\min \mathrm{costo}(t,m)$; si hay empate, $\arg\max w(t,m)$.
\end{block}
\end{columns}
\vspace{0.1cm}
{\scriptsize\color{gray} Práctica con cálculo: Llama 3.3 no supera el umbral GPQA y gana Gemma 4 26B, la arista apta más barata. Precio por millón de tokens de entrada.}
\end{frame}

% ================================================================ 5. Samuel: red de sustitución
\renewcommand{\expone}{Samuel Ramírez}
\begin{frame}{Red de sustitución: ¿con quién seguir si el modelo falla?}
\begin{columns}[T]
\column{0.55\textwidth}
\begin{itemize}
  \item Nodos: modelos aptos. Arista si son \textbf{intercambiables}: el mismo modelo en otro proveedor, o $|\Delta\mathrm{GPQA}|\le 0{,}05$.
  \item Respaldo = capas de \textbf{BFS} desde el elegido; en cada capa, el más barato.
  \item Intermediación con el algoritmo de Brandes.
\end{itemize}
\column{0.42\textwidth}
\begin{block}{Verificado con el grafo $R$ de clase}
$d(A,F)=4$, \quad $C_D(D)=0{,}6$\\
$C_B(E)=4$, \quad $C_B(D)=6{,}5$\\
$\sum\deg v = 12 = 2|E|$
\end{block}
\begin{block}{Ejemplo observado en la app}
8 aptos, 18 aristas. El elegido tiene 5 sustitutos directos ($C_D=0{,}714$).
\end{block}
\end{columns}
\[
C_D(v)=\frac{\deg v}{n-1}\qquad H(v)=\frac{1}{n-1}\sum_{u\ne v}\frac{1}{d(v,u)}\qquad C_B(v)=\sum_{\substack{s<t\\s,t\ne v}}\frac{\sigma_{st}(v)}{\sigma_{st}}
\]
{\tiny\color{gray} Fuentes: Gómez (2026), Unidad 4; Brandes (2001). Las cifras de la app dependen del catálogo disponible.}
\end{frame}

% ================================================================ 6. Samuel: inducción
\begin{frame}{Inducción: el costo oculto de una conversación larga}
\begin{columns}[c]
\column{0.55\textwidth}
Las API no tienen memoria: cada mensaje \textbf{reenvía todo el chat}.
\begin{block}{Proposición}
Tras $n$ mensajes de $t$ tokens se envían $S(n)=t\,\frac{n(n+1)}{2}$.
\end{block}
\textbf{Base:} $S(1)=t$. \ \textbf{Paso:}
{\small $S(k+1)=t\tfrac{k(k+1)}{2}+t(k+1)=t\tfrac{(k+1)(k+2)}{2}\ \blacksquare$}

\vspace{0.2cm}
\textbf{Solución:} ventana de los últimos $k$ mensajes $\Rightarrow$ total $\le n\,k\,t$, \textbf{lineal}.

\vspace{0.1cm}
{\footnotesize También por inducción: el respaldo termina en a lo más $r$ intentos.}
\column{0.42\textwidth}
\begin{tikzpicture}
\begin{axis}[width=6.1cm, height=5cm, xlabel={mensajes $n$}, ylabel={tokens enviados ($\times t$)},
  legend pos=north west, legend style={font=\scriptsize}, tick label style={font=\scriptsize}, label style={font=\small}]
\addplot[rojo, very thick, domain=0:30, samples=31] {x*(x+1)/2};
\addplot[tinta, very thick, domain=0:30, samples=31] {x<=12 ? x*(x+1)/2 : 78+12*(x-12)};
\legend{sin ventana, ventana $k=12$}
\end{axis}
\end{tikzpicture}
\end{columns}
\end{frame}

% ================================================================ 7. Manuel: resultados
\renewcommand{\expone}{Manuel Hormechea}
\begin{frame}{Resultados: 20 ejercicios de los apuntes, con APIs reales}
\begin{columns}[c]
\column{0.64\textwidth}
\renewcommand{\arraystretch}{1.3}\small
\begin{tabular}{@{}lccc@{}}
\toprule
\textbf{Estrategia} & \textbf{Aciertos} & \textbf{Cortes} & \shortstack{\textbf{Costo ref.}\\\textbf{(US\$)}} \\
\midrule
Siempre el más potente & 17/20 & 3 & 0,02906 \\
Al azar (calificados) & 20/20 & 0 & 0,01092 \\
\textbf{\color{rojo}Harness} & \textbf{20/20} & \textbf{0} & \textbf{0,00436} \\
\bottomrule
\end{tabular}
\column{0.32\textwidth}
\centering
{\fontsize{30}{34}\selectfont\bfseries\color{rojo}$-$85\,\%}\\[0.15cm]
{\footnotesize ahorro estimado de costo\\de referencia (factor 6,7)}
\end{columns}
\vspace{0.4cm}
\begin{itemize}
  \item \textbf{Observado:} 0 cortes del harness y 3 del más potente; no se registró su causa exacta.
  \item \textbf{Versión evaluada:} respaldo por costo, anterior a BFS. El salto ante 429 se prueba por separado con una API simulada.
\end{itemize}
\vspace{0.1cm}
{\scriptsize\color{gray} Fuente: resultados\_experimento.csv y .json. Una ejecución; \enquote{al azar} también obtuvo 20/20. No demuestra superioridad general en precisión.}
\end{frame}

% ================================================================ 8. Manuel: demo y conclusiones
\begin{frame}{Demostración y conclusiones}
\begin{columns}[T]
\column{0.5\textwidth}
\textbf{Demo} (\texttt{HarnessEducativo.exe}, sin cuentas)
\begin{enumerate}\small
  \item Modo \textbf{Practicar} + pregunta con cálculo.
  \item \enquote{¿Cómo lo decidió?}: red, aptos, costo.
  \item Si un modelo se satura, sigue con otro.
\end{enumerate}
\vspace{0.1cm}
\textbf{Conclusiones}
\begin{itemize}\small
  \item La teoría del curso vuelve la elección \textbf{explícita y verificable}.
  \item En la ejecución guardada: \textbf{20/20}, \textbf{0 cortes} y costo de referencia \textbf{6,7 veces menor}.
\end{itemize}
\column{0.47\textwidth}
\textbf{Limitaciones}
\begin{itemize}\small
  \item Ocho puntajes GPQA son estimados.
  \item 20 ejercicios, una sola ejecución.
  \item Red neuronal con 84 ejemplos.
  \item Selección solo al iniciar cada chat.
  \item Los cupos gratuitos cambian sin aviso.
\end{itemize}
{\small\textbf{Futuro:} calificar modelos nuevos solos y probar con estudiantes reales.}

\vspace{0.4cm}
{\Large\bfseries\color{tinta} ¡Gracias! \ {\color{rojo}¿Preguntas?}}
\end{columns}
\vspace{0.2cm}
{\footnotesize\textbf{Código y documentación:} \url{https://github.com/Jlang27/harness-educativo}}
\end{frame}

\end{document}
